\subsection*{17. Design the Role Around Interfaces and Autonomy}

Job success depends partly on what the employee receives from others and where their responsibility ends. A reviewer cannot meet a deadline if submissions arrive incomplete; a sales representative cannot promise delivery dates when operations refuses to share capacity; a faculty member cannot help learners consistently if assessment policy changes every week. A role specification should identify inputs, outputs, decision rights, escalation paths, and relationships.

The founder should ask: Which decisions can this person make independently? Which require consultation? Which require explicit approval? How does the person know whether their work is complete? Where should uncertainty go? These questions are especially important in young organizations because undocumented founder preferences can become invisible performance criteria.

Hiring someone ``senior'' to solve unclear interfaces is expensive. Seniority can improve ambiguity handling but cannot replace organizational clarity. A competent employee should have reasonable authority proportional to the responsibilities being assessed.

\subsection*{18. Hire for the Performance Horizon That Matters}

The candidate who can deliver immediately may not be the 1 who becomes strongest after 6 months. The distinction matters when skills are trainable and when the company's needs are changing quickly. A startup may require urgent coverage tomorrow; an established firm may tolerate slower ramp-up in return for greater long-term potential. Neither position is universally superior.

Define at least 3 horizons: \textbf{day-1 safety and minimum competence}, \textbf{90-day effective contribution}, and \textbf{1-year adaptability and growth}. Assess the 1st through prerequisite knowledge and bounded work samples. Assess the 2nd through representative tasks and behavioral evidence. The 3rd is harder; it may require evidence of prior learning, openness to correction, and handling unfamiliar problems. Do not equate learning potential with confidence in discussing futuristic technology.

The founder should explicitly price ramp-up and supervision. A less-experienced candidate can be a wise hire when the role is learnable and mentorship exists. The same choice can be unfair to the person and dangerous to customers if the company expects independent mastery from day 1.

\subsection*{19. Role Specifications Are Hypotheses That Must Be Revised}

At a new venture, the founder may genuinely not know which competencies predict success. A 1st role profile is therefore a structured hypothesis. It should be written down before recruitment, tested against actual work, and revised as evidence accumulates. What must not change casually is the scoring standard \textbf{during a particular candidate competition}, because shifting criteria after seeing applicants destroys comparability.

After the hire has performed real work, ask which predicted competencies mattered, which selection exercises were misleading, and what unexpected tasks consumed time. Perhaps technical depth was already common among applicants, while response reliability differentiated performance. Perhaps the supposedly critical tool was replaced within weeks, while diagnostic reasoning remained important.

This feedback is how job analysis becomes a living institutional asset. A business that never revisits specifications will eventually select for a past version of the work. A business that changes them after each attractive candidate will never learn what it is selecting for.

\subsection*{20. Draft an Honest Job Advertisement After the Internal Work Is Clear}

The advertisement is the candidate-facing expression of the success profile. It should state actual responsibilities, realistic compensation information where appropriate or required, working arrangements, required schedule, necessary credentials, expected experience, assessment steps, and what would make someone successful. The language should be understandable to someone outside the company. Describe the company's stage truthfully: a founder-led new venture should not imply the security, hierarchy, or institutional prestige of an established employer.

Distinguish what is essential from what is preferred. Avoid adding a degree requirement because comparable firms include 1, or demanding 5 years in a tool that has existed for fewer years. Such filters may reduce applicant volume but do not necessarily improve quality. When talent is scarce, overrestriction can be more expensive than additional screening.

A well-written advertisement also contains a \textbf{realistic job preview}: attractive features and difficult features. Candidates should understand ambiguity, intensity, autonomy, repetitive work, customer exposure, and boundaries. An applicant who self-selects out after reading honest requirements may save both sides a costly mismatch.

\subsection*{21. Design the Full System Before Opening Applications}

Before posting, decide: who owns the role; the compensation ceiling; the permissible employment or contracting relationship; the sourcing channels; essential qualifications; the assessment sequence; the scoring rubric; who evaluates; target response times; how references and data will be handled; and who has authority to make an offer. These choices belong together because a process can fail even when each isolated assessment appears sensible.

For example, if 250 applicants are expected and a 2-hour work sample is required from every applicant, the company imposes 500 candidate-hours before any meaningful screening. If 2 evaluators spend 30 minutes on each submission, the business incurs 250 reviewer-hours. That may be economically absurd for a low-volume or modestly paid role. A shorter, fair preliminary screen followed by a deeper finalist assessment can produce more information per unit of burden.

Good process design also anticipates accommodations, technical failures, candidate questions, and evaluator absence. Hiring is an operational service delivered to candidates. Reliability matters on both sides.

\subsection*{22. Treat the Absence of a Role as a Decision with an Opportunity Cost}

Founders sometimes feel compelled to hire because revenue is increasing, investors expect headcount, or peers have begun building departments. But the right timing depends on incremental value. An additional hire might unlock more delivery, reduce founder bottlenecks, improve quality, or preserve essential continuity. They might also create supervision overhead and fixed cost before demand exists.

Compare the \textbf{expected contribution created or protected by the role} with total compensation, tools, training, management time, error risk, and contractual commitments. Then compare alternatives: fractional specialists, task redesign, temporary support, or delaying an offering. If a person is hired, the job should be capable of producing valuable, observable work: not merely providing the founder with the feeling of organizational progress.

This economic question does not reduce people to line items. It protects them from being recruited into an unsustainable role and protects the institution from making commitments it cannot honor. The most humane offer is often 1 backed by genuine work, realistic management capacity, and sufficient cash.
